\documentclass[10pt,a4paper]{amsart}
\usepackage[margin=19mm]{geometry}
\usepackage{amsmath,amssymb,mathtools}
\usepackage[scr=rsfs,cal=euler]{mathalfa}
\usepackage{xfrac}
\usepackage[unicode,hidelinks,bookmarks=false]{hyperref}
\newcommand{\ie}{i.e.\ }
\newcommand{\cf}{cf.\ }
\newcommand{\resp}{respectively\ }
\newcommand{\Subsection}{\subsection}
\providecommand{\nobreakdash}{\nobreak\hbox{-}}
\setlength{\emergencystretch}{2em}
\multlinegap0pt
\usepackage{bm}
\usepackage{bbm}
\usepackage{dsfont}
\usepackage{xspace}
\usepackage{cancel}
\usepackage{mathtools}
\usepackage{cleveref}
\usepackage{amsthm}
\makeatletter
\@ifundefined{theorem}{\newtheorem{theorem}{Theorem}}{}
\@ifundefined{definition}{\newtheorem{definition}[theorem]{Definition}}{}
\@ifundefined{remark}{\newtheorem{remark}[theorem]{Remark}}{}
\makeatother
\newcommand{\supp}{{\mathsf{supp }\hspace{.075cm}}}
\title[Stochastically Evolving Graphs via Edit Semigroups]{Stochastically Evolving Graphs via Edit Semigroups}
\author{Fan Chung, Sawyer Jack Robertson}
\date{Source: arXiv:2509.19678v1 (CC BY 4.0)}
\begin{document}
\maketitle

\noindent\emph{Source and licence.} Stochastically Evolving Graphs via Edit Semigroups. Version arXiv:2509.19678v1 (CC BY 4.0), distributed under the Creative Commons Attribution 4.0 International licence, as recorded in the arXiv licence metadata for this exact version and shown on the abstract page https://arxiv.org/abs/2509.19678v1. Full rights notice: licenses/arxiv-2509.19678-CC-BY-4.0.txt.\par\smallskip

\noindent\textit{Editorial scope.} Counted unit: the Theorem whose source label is \texttt{thm:orthonormality} in the article (source0.tex lines 685--691, source label \texttt{thm:orthonormality}), delivered together with the article's own complete original proof of it (source0.tex lines 693--748, one \texttt{proof} environment of 56 lines). No mathematics is written, repaired or re-derived: statement and proof are the article's own text line for line. Required context retained uncounted: def:graph-edit-semigroup (lines 297--303); eq:edit-introducton (lines 299--301); eq:semi-rw (lines 375--377); thm:eigenvalues-state-graph (lines 382--396); rmk:multiplicities (lines 398--400); eq:eigenvector-defn (lines 637--639). Every definition, display and label of those lines is retained unchanged. Omitted from this package and not redistributed: 1--296, 302--374, 378--381, 397--397, 401--636, 640--684, 692--692, 749--1200. Nothing inside a retained excerpt is omitted or altered: every retained body is delivered exactly as the article's lines are written, so no omission receipt of any kind accompanies this package. Citations: the retained excerpts carry no citation command at all, so no bibliography entry of the article is transcribed and no reference is redistributed. Reference adaptation: none was needed and none was made; every reference inside the delivered unit resolves inside it, so no unresolved reference or citation is produced. Printed slips kept literal: no printed word, symbol, hypothesis or number of the counted statement or of its proof was altered. Layout and class: the article's own class is \texttt{article} with packages including inputenc, lmodern, fontenc, babel, ifpdf, geometry, xcolor, hyperref, natbib, subcaption, caption, graphicx, amsmath, amsthm; the wrapper is the lane's pinned \texttt{amsart} editorial wrapper at 10pt on A4, which restates the theorem-like environments the retained text needs and adds the article's own macro definitions that the retained text needs (\texttt{\textbackslash supp}), with extra wrapper packages (bm, bbm, dsfont, xspace, cancel, mathtools, cleveref). The delivered numbering is the unit's own. Assets: no retained span contains a figure environment, a graphics-inclusion command, a PDF-page inclusion, a table or a TikZ picture; no image file is copied into this package, the package declares no asset dependency and no image file is redistributed with it. Licence: version arXiv:2509.19678v1 of this article is distributed under the Creative Commons Attribution 4.0 International licence, as recorded in the arXiv licence metadata for this exact version and shown on the abstract page https://arxiv.org/abs/2509.19678v1. A full rights notice is delivered at licenses/arxiv-2509.19678-CC-BY-4.0.txt. Characters: every retained excerpt is pure ASCII, so the delivered engine, XeLaTeX with its default Latin Modern text font, reports no missing character.
\begin{definition}[Graph edit semigroup]\label{def:graph-edit-semigroup}
    Let $\mathcal{H}=(\mathcal{V},\mathcal{E})$ be a fixed host graph and let $E\subseteq\mathcal{E}$ be a subset of edges. For each $e\in \mathcal{E}$, let the symbol ${e^+}$ denote the operation of ``add edge $e$ to $E$'' and let ${e^-}$ denote the operation of ``remove edge $e$ from $E$''. We call these {\normalfont simple edits}. We write, for $e_i\in {\mathcal{E}}$ and $\sigma_i\in\{\pm\}$,
        \begin{align}\label{eq:edit-introducton}
            e_k^{\sigma_k} e_{k-1}^{\sigma_{k-1}} \cdots e_{1}^{\sigma_1} E
        \end{align}
    to represent the graph obtained by {\normalfont first} performing the edit $e_{1}^{\sigma_1}$, then so on, until performing the edit $e_k^{\sigma_k}$. If $x, y$ are sequences of edits such as those applied to $E$ in~\cref{eq:edit-introducton}, we write $x = y$ if $xE = yE$ for each $E\subseteq\mathcal{E}$. The semigroup $\mathcal{S} = \mathcal{S}(\mathcal{H})$ generated by the set of simple edits is called the {\normalfont graph edit semigroup}. By convention, $\mathcal{S}$ contains an identity edit denoted $\widehat{0}$.
\end{definition}

    \begin{align}\label{eq:semi-rw}
        P(x, y) = \sum_{\substack{z\in S\\zx=y}} w_z, \quad x, y\in S.
    \end{align}

\begin{theorem}\label{thm:eigenvalues-state-graph}
    Let $\mathcal{S} = \mathcal{S}(\mathcal{H})$ denote the graph edit semigroup corresponding to the host graph $\mathcal{H}$ as in~\cref{def:graph-edit-semigroup}. Let $(w_x)_{x\in \mathcal{S}}$ be a probability distribution on $\mathcal{S}$, and let $\mathsf{P}$ be the Markov transition probability matrix for the random walk on the chambers of $\mathcal{S}$ given in~\cref{eq:semi-rw}. Then $\mathsf{P}$ is diagonalizable and has an eigenvalue with multiplicity one for each subset $T\subseteq {\mathcal{E}}$ given by
        \begin{align}\label{eq:eigenvalue-01}
            \lambda_T = \sum_{\substack{y\in \mathcal{S} \\\supp{y} \subseteq T}} w_y.
        \end{align}
    In the special case where $(w_x)_{x\in S}$ is nonzero exactly on the set of simple edits with edit probabilities determined by $\mathbf{p} = (p_e)_{e\in\mathcal{E}}$ such that $0<p_e<1$ and which satisfies
        \begin{align*}
            w_{e^+} = 1 - w_{e^-} = \frac{p_e}{{m}},
        \end{align*}
    then the semigroup random walk on the chambers $\mathcal{C}$ of $\mathcal{S}$ has the same transition probability matrix as the simple edit process with edge probabilities $\mathbf{p}$. Moreover, its eigenvalues are indexed by subsets $T\subseteq\mathcal{E}$ and for each such subset we have
        \begin{align}\label{eq:lambda-T}
            \lambda_T = \frac{|T|}{{m}},
        \end{align}
    independent of $(p_e)_{e\in\mathcal{E}}$.
\end{theorem}

\begin{remark}\label{rmk:multiplicities}
    We remark that the eigenvalue $\lambda_T$ in~\cref{thm:eigenvalues-state-graph} that is \textit{indexed by} $T\subseteq\mathcal{E}$ occurs with multiplicity $m_T=1$, but distinct $T$ may give rise to the same value $\lambda\in\mathbb{R}$. In the case of the simple edit process, each eigenvalue of the form $\frac{k}{{m}}$ for some $0\leq k\leq m$ occurs with multiplicity ${{m}\choose k}$.
\end{remark}

        \begin{align}\label{eq:eigenvector-defn}
            \phi_T(E) &= (-1)^{|\mathcal{E}\setminus(E\cup T)|}\prod_{e\in T\cap E}p_e \prod_{e\in T\setminus E}(1-p_e),\quad E\subseteq \mathcal{E}.
        \end{align}

\begin{theorem}\label{thm:orthonormality}
    Let $\mathcal{H}=(\mathcal{V}, \mathcal{E})$ be a given host graph and let $\mathsf{P}$ (resp. $\pi$) denote the transition probability matrix (resp. stationary distribution) of the simple edit process with edge probabilities $\mathbf{p} = (p_e)_{e\in\mathcal{E}}$. For each $T\subseteq\mathcal{E}$, let $\phi_T$ denote the correspondingly indexed eigenvector of $\mathsf{P}$ as in~\cref{eq:eigenvector-defn}. Let $\psi_T\in\mathbb{R}^{2^\mathcal{E}}$ be given by the formula
        \begin{align*}
            \psi_T(E) &=\phi_T(E)  \frac{\prod_{e\notin T} \sqrt{p_e(1-p_e)}}{\sqrt{\pi(E)}},\quad E\subseteq\mathcal{E}.
        \end{align*}
    Then $\{\psi_T\}_{T\subseteq \mathcal{E}}$ form an orthonormal system of (left) eigenvectors for the symmetric matrix $\mathsf{Q} = \Pi^{1/2}\mathsf{P}\Pi^{-1/2}$, where $\Pi = \mathrm{diag}(\pi(E))_{E\subseteq\mathcal{E}}$.
\end{theorem}

\begin{proof}
    Note that, up to the scaling factor $\prod_{e\notin T} \sqrt{p_e(1-p_e)}$ which is uniform in $E\subseteq{\mathcal E}$, it holds $\psi_T = \phi_T \Pi^{-1/2}$. Thus $\psi_T \mathsf{Q}= \phi_T\mathsf{P}\Pi^{-1/2} = \lambda_T \psi_T$, and it follows that $\{\psi_T\}_{T\subseteq \mathcal{E}}$ forms a complete collection of eigenvectors of $\mathcal{Q}$. 

    The symmetry of $\mathsf{Q}$ guarantees orthogonality of $\psi_T, \psi_{T'}$ provided $\lambda_T\neq\lambda_{T'}$; however the spectrum of $\mathsf{P}$ will in general contain many repeated eigenvalues (see~\cref{rmk:multiplicities}) and thus we check orthogonality in general as follows. For a given subset $E\subseteq\mathcal{E}$, let $\mathbf{1}_E\in\mathbb{R}^{\mathcal{E}}$ denote the indicator vector of $E$. Then for any $T\subseteq\mathcal{E}$ fixed, the vector $\phi_T$ may be recast as a product over edges in the form
        \begin{align*}
            \phi_T(E)=\prod_{e\in\mathcal E}g_{e,T}(\mathbf{1}_E(e)),\quad E\subseteq\mathcal{E},
        \end{align*}
    where $g_{e,T}:\{0, 1\}\rightarrow\mathbb{R}$ is given by
        \begin{align*}
            g_{e,T}(x)=
            \begin{cases}
            p_e,&x=1\text{ and }e\in T,\\
            1-p_e,&x=0\text{ and }e\in T,\\
            1,&x=1\text{ and }e\notin T,\\
            -1,&x=0\text{ and }e\notin T.
            \end{cases}.
        \end{align*}
    Hence, for $S, T\subseteq\mathcal{E}$ fixed, it holds
        \begin{align*}
            \langle\psi_T,\psi_S\rangle
            &=\sum_{E\subseteq\mathcal E}\psi_T(E)\,\psi_S(E)\\
            &=\prod_{e\notin T}\sqrt{p_e(1-p_e)}
            \;\prod_{e\notin S}\sqrt{p_e(1-p_e)}\\
            &\quad\times \sum_{X\in\{0,1\}^{\mathcal E}}
            \frac{\prod_{e}g_{e,T}(X(e))\,g_{e,S}(X(e))}{\pi(E)},
        \end{align*}
    where we identify subsets with their indicator vectors. Focusing on the last term for a moment, we have the factorization:
        \begin{align*}
            \sum_{X\in\{0,1\}^{\mathcal E}}&\frac{\prod_{e\in\mathcal{E}}g_{e,T}(X(e))\,g_{e,S}(X(e))}{\pi(E)}\\
            &=\sum_{X\in\{0,1\}^{\mathcal E}} \prod_{e\in\mathcal{E}} \frac{g_{e,T}(X(e)) g_{e,S}(X(e))}{p_e^{X(e)}(1-p_e)^{1-X(e)}}\\
            &=\prod_{e\in\mathcal E}
            \Bigl(\sum_{x=0}^1\frac{g_{e,T}(x)\,g_{e,S}(x)}{p_e^x(1-p_e)^{1-x}}\Bigr).
        \end{align*}
    This can be seen by expanding $\sum_{X\in\{0,1\}^{\mathcal E}}(\cdot)$ into $|\mathcal{E}|$ nested sums and then rearranging. If $T\neq S$, choose an edge $e^\ast$ on which they differ; assuming without loss of generality that $e^\ast\in T$ and $e^\ast\notin S$, one has
        \begin{align*}
            \sum_{x=0}^1\frac{g_{e^\ast,T}(x)\,g_{e^\ast,S}(x)}{p_{e^\ast}^x(1-p_{e^\ast})^{1-x}}
            &=-\frac{1-p_e}{1-p_e} + \frac{p_e}{p_e} = 0.
        \end{align*}
    so $\langle\psi_T,\psi_S\rangle=0$. On the other hand, if $T=S$, then for each $e$ one checks
        \begin{align*}
        \sum_{x=0}^1\frac{g_{e,T}(x)^2}{p_e^x(1-p_e)^{1-x}}
        &=
        \begin{cases}
        \frac{(1-p_e)^2}{1-p_e}+\frac{p_e^2}{p_e}=1,&e\in T,\\
        \frac{1}{1-p_e}+\frac{1}{p_e}=\frac{1}{p_e(1-p_e)},&e\notin T.
        \end{cases}
        \end{align*}
    Therefore
        \begin{align*}
        \langle\psi_T,\psi_T\rangle
        &=\prod_{e\notin T}p_e(1-p_e)
        \;\prod_{e\notin T}\frac1{p_e(1-p_e)}
        =1,
        \end{align*}
    which completes the proof.
\end{proof}

\end{document}
